\documentclass[12 pt, letterpaper]{hmcpset}
\usepackage{hw-preamble}

\name{Jayden Thadani}
\class{MATH 171 --- Abstract Algebra I}
\assignment{Homework 13}
\duedate{5th April 2024}

\begin{document}

\begin{problem}[Dummit \& Foote \S7.1 \#14]
	Let $x$  be a nilpotent element of the commutative ring $R$ 
	\begin{enumerate}
		\item Prove that $x$ is either zero or a zero divisor.
		\item Prove that $rx$ is nilpotent for all $r \in R$.
		\item Prove that $1 + x$ is a unit in $R$.
		\item Deduce that the sum of a nilpotent element and a unit is a unit.
	\end{enumerate}
\end{problem}

\begin{solution}
	\begin{enumerate}
		\item Suppose $x$ is nonzero. Then there exists some least $n \in \mathbb{N}$ such that $x^{n} = 0$ with $n > 1$. Thus, $x x^{n - 1} = 0$ and $x^{n - 1}$ is nonzero (since $n$ is the least positive integer that gives $x^{n} = 0$). That is, $x$ is a zero divisor.
		\item Suppose $x^{n} = 0$. Then, for any $r \in R$, we have
			\[
				\begin{split}
					(rx)^{n} & = \underbrace{(rx) \dots (rx)}_{n \text{ times}} \\
						 & = r^{n} x^{n} \\
						 & = r^{n} 0 \\
						 & = 0.
				\end{split}
			\]
			where the second step follows from the commutativity of $R$. Thus, $rx$ is nilpotent.
		\item Suppose $x^{n} = 0$. Note that the polynomial $1 - y^{n} \in \mathbb{Z}[x]$ has root $y = - 1$ if $n$ is even and the polynomial $1 + y^{n}$ has root $y = -1$ if $n$ is odd. Specifically, if $n$ is even, we can write
			\[
				\begin{split}
					1 - y^{n} & = (1 - y)(y^{n - 1} + y^{n - 2} + \dots + y + 1) \\
						  & = (1 - y)(1 + y)(1 + y^{2} + \dots + y^{n - 2})
				\end{split}
			\]
			and
			\[
				1 + y^{n} = (1 + y)(y^{n - 1} - y^{n - 2} + \dots + y^{2} - y + 1).
			\]
			Since multiplication distributes over addition just as well in $R$ as it does in $\mathbb{Z}[x]$, we can write
			\[
				\begin{split}
					(1 + x)(1 - x)(1 + x^{2} + \dots + x^{n - 2}) & = 1 - x^{n} \\
										      & = 1 - 0 \\
										      & = 1
				\end{split}
			\]
			for $n$ even and
			\[
				\begin{split}
					(1 + x)(x^{n - 1} - x^{n - 2} + \dots + x^{2} - x + 1) & = 1 + x^{n} \\
											       & = 1 + 0 \\
											       & = 1
				\end{split}
			\]
			for $n$ odd. Thus, we have that $1 + x$ is a unit in both cases.
		\item We can use a similar proof to show that for any unit $u \in R$, $u + x$ is a unit. 
			
			First note that for any $n \in \mathbb{N}$, $u^{n}$ is a unit --- just consider $u^{n} (u^{-1})^{n} = 1$. Also note that for any unit $u$, if $a, b \in R$ give $ab = u$, then $a, b$ must be units too. We can see this by first noticing if one is a unit, then the other must be. Without loss of generality, suppose $a$ is a unit. Then $b = a^{-1} u$ must be a unit since $u^{-1} a a^{-1} u = 1$. Notice that $a$ must be a unit since $a (b u^{-1}) = 1$, and thus both must be units. 
			
			Then by the same factorisation trick, for $n$ such that $x^{n} = 0$, we can write $u^{n} + x^{n}$ or $u^{n} - x^{n}$ (depending on whether $n$ is odd or even) as $(u + x) p(x)$ for some $p \in \mathbb{Z}[x]$. We can use the canonical homomorphism of $\mathbb{Z}$ into every ring to replace the coefficients $a \in \mathbb{Z}$. That is, use $\varphi : \mathbb{Z} \to R$ by $1 \mapsto 1_{R}$ to write $a \mapsto \underbrace{1 + \dots + 1}_{a \text{ times}}$. 
			
			Since distributivity still works, we have for all $n \in \mathbb{N}$
			\[
				(u + x) r = u^{n}
			\]
			for some $r \in R$. Then, by our previous results that $u^{n}$ is a unit and that two elements that multiply to form a unit must be a unit, we have that $u + x$ is a unit.
	\end{enumerate}
\end{solution}

\pagebreak

\begin{problem}[Dummit \& Foote \S7.2 \#10]
	Consider the following elements of the integral group ring $\mathbb{Z} S_{3}$ :
	\[
		\alpha = 3 (1 \, 2) - 5 (2 \, 3) + 14 (1 \, 2 \, 3) \quad \beta = 6 (1) + 2 (2 \, 3) - 7 (1 \, 3 \, 2)
	\]
	(where $(1)$ is the identity of $S_{3}$) . Compute the following elements:
	\begin{enumerate}
		\item $\alpha + \beta$
		\item $2 \alpha - 3 \beta$
		\item $\alpha \beta$
		\item $\beta \alpha$
		\item $\alpha^{2}$
	\end{enumerate}
\end{problem}

\begin{solution}
	\begin{enumerate}
		\item
			\[
				\boxed{
					\alpha + \beta = 6 (1) + 3 (1 \, 2) - 3 (2 \, 3) + 14 (1 \, 2 \, 3) - 7 (1 \, 3 \, 2).
				}
			\] 

			It follows by the standard axioms of distributivity and commutativity of addition that
			\[
				\begin{split}
					\alpha + \beta & = 3 (1 \, 2) - 5 (2 \, 3) + 14 (1 \, 2 \, 3) + 6 (1) + 2 (2 \, 3) - 7 (1 \, 3 \, 2) \\
						       & = 6 (1) + 3 (1 \, 2) + (-5 + 2)(2 \, 3) + 14 (1 \, 2 \, 3) - 7 (1 \, 3 \, 2) \\
						       & = 6 (1) + 3(1 \, 2) - 3(2 \, 3) + 7(1 \, 2 \, 3).
				\end{split}
			\]
		\item
			\[
				\boxed{
					2 \alpha - 3 \beta = -18 (1) + 6 (1 \, 2) - 16 (2 \, 3) + 49 (1 \, 2 \, 3).
				}
			\] 

			Again, using distributivity and commutativity of addition, we have
			\[
				\begin{split}
					2 \alpha - 3 \beta & = 2(3 (1 \, 2) - 5 (2 \, 3) + 14 (1 \, 2 \, 3)) - 3(6 (1) + 2 (2 \, 3) - 7 (1 \, 3 \, 2)) \\
							   & = 6 (1 \, 2) - 10(2 \, 3) + 28 (1 \, 2 \, 3)- 18 (1) - 6(2 \, 3) + 21(1 \, 3 \, 2) \\
							   & = -18 (1) + 6 (1 \, 2) + (-10 - 6)(2 \, 3) + 28 (1 \, 2 \, 3) + 21 (1 \, 3 \, 2) \\
							   & = -18 (1) + 6(1 \, 2) -16 (2 \, 3) + 28 (1 \, 2 \, 3) + 21 (1 \, 3 \, 2).
				\end{split}
			\]
		\item
			\[
				\boxed{
					\alpha \beta = -108 (1) + 81 (1 \, 2) - 30 (2 \, 3) - 21 (3 \, 1) + 90 (1 \, 2 \, 3).
				}
			\] 

			Here we use distributivity and commutativity of addition as usual, but we also need to use the fact that elements of $\mathbb{Z}$ commute with elements of $S_{3}$ under multiplication for the second step.
			\[
				\begin{split}
					\alpha \beta & = (3(1 \, 2) - 5 (2 \, 3) + 14(1 \, 2 \, 3))(6(1) + 2(2 \, 3) - 7(1 \, 3 \, 2)) \\
						     & = (3 \times 6) (1 \, 2) (1) - (5 \times 6) (2 \, 3) (1) + (14 \times 6)(1 \, 2 \, 3)(1) \\
						     & + (3 \times 2)(1 \, 2)(2 \, 3) + (-5 \times 2)(2 \, 3)(2 \, 3) + (14 \times 2)(1 \, 2 \, 3)(2 \, 3) \\
						     & + (3 \times(-7))(1 \, 2)(1 \, 3 \, 2) + ((-5) \times (-7))(2 \, 3)(1 \, 3 \, 2) + (14 \times (-7))(1 \, 2 \, 3)(1 \, 3 \, 2) \\
						     & = 18 (1 \, 2) - 30 (2 \, 3) + 84 (1 \, 2 \, 3) + 6 (1 \, 2 \, 3) - 10(1) \\
						     & + 28 (2 \, 1) - 21 (1 \, 3) + 35 (2 \, 1) - 98 (1) \\
						     & = (-10 - 98)(1) + (18 + 28 + 35)(1 \, 2) - 30 (2 \, 3) - 21 (1 \, 3) + (84 + 6)(1 \, 2 \, 3) \\
						     & = -108 (1) + 81 (1 \, 2) - 30 (2 \, 3) - 21 (3 \, 1) + 90 (1 \, 2 \, 3).
				\end{split}
			\]
		\item
			\[
				\boxed{
					\beta \alpha = -108 (1) + 18 (1 \, 2) - 51 (2 \, 3) + 63 (3 \, 1) + 84 (1 \, 2 \, 3) + 6 (1 \, 3 \, 2).
				}
			\]

			Like for $\alpha \beta$ this follows from distributivity, commutativity of addition and $\mathbb{Z}$ commuting with $S_{3}$.
			\[
				\begin{split}
					\beta \alpha & = (6(1) + 2(2 \, 3) - 7(1 \, 3 \, 2))(3(1 \, 2) - 5 (2 \, 3) + 14(1 \, 2 \, 3)) \\
						     & = (6 \times 3) (1) (1 \, 2) - (6 \times 5) (1) (2 \, 3) + (6 \times 14) (1) (1 \, 2 \, 3) \\
						     & + (2 \times 3) (2 \, 3) (1 \, 2) + (2 \times (-5)) (2 \, 3) (2 \, 3) + (2 \times 14) (2 \, 3) (1 \, 2 \, 3) \\
						     & + ((-7) \times 3) (1 \, 3 \, 2) (1 \, 2) + ((-7) \times (-5)) (1 \, 3 \, 2) (2 \, 3) + ((-7) \times 14) (1 \, 3 \, 2) (1 \, 2 \, 3) \\
						     & = 18 (1 \, 2) - 30 (2 \, 3) + 84 (1 \, 2 \, 3) + 6 (1 \, 3 \, 2) - 10(1) \\
						     & + 28 (1 \, 3) - 21 (2 \, 3) + 35 (1 \, 3) - 98 (1) \\
						     & = (-10 - 98)(1) + 18 (1 \, 2) + (-30 - 21) (2 \, 3) + (28 + 35) (1 \, 3) + 84 (1 \, 2 \, 3) \\
						     & = -108 (1) + 18 (1 \, 2) - 51 (2 \, 3) + 63 (3 \, 1) + 84 (1 \, 2 \, 3) + 6 (1 \, 3 \, 2).
				\end{split}
			\]
		\item
			\[
				\boxed{
					\alpha^{2} = -16 (1) - 70 (1 \, 2) + 42 (2 \, 3) - 28 (3 \, 1) -15 (1 \, 2 \, 3) + 181 (1 \, 3 \, 2).
				}
			\]

			For $\alpha^{2}$ we can use the identity that $(a + b + c)^{2} = a^{2} + b^{2} + c^{2} + ab + bc + ca + ba + cb + ac$ which only requires distributivity, and thus, works in any ring.
			\[
				\begin{split}
					\alpha^{2} & = (3 (1 \, 2) - 5 (2 \, 3) + 14 (1 \, 2 \, 3))^{2} \\
						   & = 3^{2} (1 \, 2)^{2} - 5^{2} (2 \, 3)^{2} + 14^{2} (1 \, 2 \, 3)^{2} \\ 
						   & + (3 \times (-5))(1 \, 2)(2 \, 3) + ((-5) \times 14) (2 \, 3) (1 \, 2 \, 3) + (14 \times 3) (1 \, 3 \, 2)(1 \, 2) \\
						   & + ((-5) \times 3)(2 \, 3)(1 \, 2) + (14 \times (-5))(1 \, 2 \, 3)(2 \, 3) + (3 \times 14)(1 \, 2)(1 \, 3 \, 2) \\
						   & = 9 (1) - 25 (1) + 196 (1 \, 3 \, 2) - 15 (1 \, 2 \, 3) - 70 (1 \, 3) + 42 (2 \, 3) \\
						   & - 15 (1 \, 3 \, 2) - 70 (1 \, 2) + 42 (1 \, 3) \\
						   & = (9 - 25) (1) - 70 (1 \, 2) + 42(2 \, 3) + (42 - 70) (1 \, 3) + (196 - 15)(1 \, 3 \, 2) - 15 (1 \, 2 \, 3) \\
						   & = -16 (1) - 70 (1 \, 2) + 42 (2 \, 3) - 28 (3 \, 1) -15 (1 \, 2 \, 3) + 181 (1 \, 3 \, 2).
				\end{split}
			\]
	\end{enumerate}
\end{solution}

\pagebreak

\begin{problem}[Dummit \& Foote \S7.4 \#7]
	Let $R$ be a commutative ring with $1$. Prove that the principal ideal generated by $x$ in the polynomial ring $R[x]$ is a prime ideal if and only if $R$ is an integral domain. Prove that $(x)$ is a maximal ideal if and only if $R$ is a field.
\end{problem}

\begin{solution}
	Note that $(x)$ being a prime ideal in $R[x]$ is equivalent to $R[x]/(x)$ being an integral domain and $(x)$ being maximal in $R[x]$ is equivalent to $R[x]/(x)$ being a field. We claim that $R[x]/(x) \cong R$ and thus, $(x)$ being a prime ideal is equivalent to $R$ being an integral domain and $(x)$ being a maximal ideal is equivalent to $R$ being a field. 

	
	This proof rests on the following lemmas ---
	\begin{lemma}
		If $R \cong S$ (commutative rings with identity) and $S$ is an integral domain, then $R$ is an integral domain.
	\end{lemma}
	\begin{proof}
		Let $\varphi$ be an isomorphism $\varphi : R \to S$. Suppose $a, b \in R \setminus \{ 0 \}$. Then $\varphi(a)$ and $\varphi(b)$ are nonzero, since isomorphisms must be injective. Since $S$ is an integral domain, $\varphi(a)\varphi(b) = \varphi(ab)$ is nonzero, and thus, $ab$ is nonzero (if it were $0$, then $\varphi(ab) = 0$). Thus $R$ is an integral domain
	\end{proof}
	\begin{lemma}
		If $R \cong S$ (commutative rings with identity) and $S$ is a field, then $R$ is a field.
	\end{lemma}
	\begin{proof}
		Let $\varphi$ be an isomorphism $\varphi : R \to S$. Suppose $a \in R \setminus \{ 0 \}$. Then $\varphi(a)$ is also nonzero, and since $S$ is a field, it is a unit. Thus there exists some $\varphi(a)^{-1} \in S$. Since $\varphi$ is an isomorphism, there also exists some $\varphi^{-1}(\varphi(a)^{-1}) \in R$. 

		Note that
		\[
			\begin{split}
				\varphi^{-1}(\varphi(a)^{-1}) a & = \varphi^{-1}(\varphi(a)^{-1}) \varphi^{-1}(\varphi(a)) \\
								& = \varphi^{-1}(\varphi(a)^{-1} \varphi(a)) \\
								& = \varphi^{-1}(1_{S}).
			\end{split}
		\]
		We claim that $\varphi^{-1}(1_{S}) = 1_{R}$. Notice that for any $r \in R$,
		\[
			\begin{split}
				\varphi^{-1}(1_{S}) r & = \varphi^{-1}(1_{S}) \varphi^{-1}(\varphi(r)) \\
						      & = \varphi^{-1}(1_{S} \varphi(r)) \\
						      & = \varphi^{-1}(\varphi(r)) \\
						      & = r.
			\end{split}
		\]
		Thus, $\varphi^{-1}(1_{S})$ is the unique multiplicative identity in $R$ and $\varphi^{-1}(\varphi(a)^{-1}) a = 1_{R}$, giving us that $a$ is a unit. Thus, any nonzero element of $R$ is a unit, and $R$ is a field.
	\end{proof}
	We will show that $R[x]/(x) \cong R$ by the first isomorphism theorem. Specifically, consider the ring homomorphism
	\[
		\begin{split}
			\varphi \colon & R[x] \to R \\
				       & p \mapsto p(0)
		\end{split}
	\]
	(we've shown in class that evaluation is a ring homomorphism). 

	
	We claim that $\varphi$ is surjective and $\ker \varphi = (x)$. We can see that $\varphi$ is surjective, since for any $a \in R$, the constant polynomial $p = a$ gives $\varphi(p) = a$.  Also note that for any polynomial $p = a_{n} x^{n} + \dots + a_{1} x + a_{0}$, we have $p(0) = a_{0}$. Thus, the kernel of $\varphi$ is just those polynomials $p$ with $a_{0} = 0$, which is just $(x)$. 

	
	Then by the first ring isomorphism theorem, we have $R[x]/\ker \varphi \cong \operatorname{im} \varphi$, which is $R[x]/(x) \cong R$ as desired. 

	
\end{solution}

\pagebreak

\begin{problem}[Dummit \& Foote \S7.4 \#8]
	Let $R$ be an integral domain. Prove that $(a) = (b)$ for some elements $a, b \in R$ if and only if $a = ub$ for some unit $u$ of $R$.
\end{problem}

\begin{solution}	
	Suppose $(a) = (b)$, for nonzero $a, b$. Then since $a \in (a)$ and $b \in (b)$, we have that. $a \in (b)$ and $b \in (a)$. Thus, $a = br_{1}$ and $b = a r_{2}$ for some $r_{1}, r_{2} \in R$. Then, by substituting $ar_{2}$ for $b$ we have $a = a r_{2} r_{1}$. 

	If $a = 0$ then $(a) = \{ 0 \}$ and thus, if $(b) = (a)$, $b \in \{ 0 \}$ and $b = 0$ as well. Then $a = 1 b$. as desired. 
	
	If $a$ is nonzero, then since $R$ is an integral domain we can cancel to get $r_{1} r_{2} = 1$. Thus, $r_{1} = u$ is the unit such that $a = ub$. 

	Suppose $a = ub$ for some unit $u$. Notice that any $r \in (a)$ must be of the form $ar_{1}$ for some $r_{1} \in R$. But $ar_{1} = b u r_{1}$ and thus, $r = u r_{1} b \in (b)$. That is $(a) \subset (b)$. Since $b = u^{-1} a$, by the same proof we also have $(b) \subset (a)$. Thus $(a) = (b)$, completing the proof.
\end{solution}

\pagebreak

\begin{problem}[Dummit \& Foote \S7.4 \#11]
	Assume that $R$ is commutative. Let $I$ and $J$ be ideals of $R$ and assume $P$ is a prime ideal of $R$ that contains $IJ$ (for example, if $P$ contains $I \cap J$). Prove either $I$ or $J$ is contained in $P$.
\end{problem}

\begin{solution}
	If $P$ contains $J$, we're done. If $P$ doesn't contain $J$, then there must be some specific $j \in J$ such that $j \notin P$. But for all $i \in I$, $ij \in IJ$ and thus $ij \in P$. Since $j \notin P$, and $P$ is a prime ideal, we must have $i \in P$. That is, for every $i \in I$ we have $i \in P$ --- we have $I \subset P$.
\end{solution}

\end{document}
